\documentclass[12pt,a4paper]{article}

% Paquetes
\usepackage[spanish]{babel}
\usepackage[utf8]{inputenc}
\usepackage{amsmath, amssymb}
\usepackage{geometry}
\usepackage{graphicx}
\usepackage{float}

\geometry{margin=2.5cm}

% Título
\title{Método de Euler y Análisis de Error}
\author{}
\date{}

\begin{document}

\maketitle

\section{Introducción}

El método de Euler es uno de los métodos numéricos más simples para aproximar la solución de una ecuación diferencial ordinaria (EDO) de primer orden de la forma:

\begin{equation}
\frac{dy}{dx} = f(x,y), \quad y(x_0) = y_0
\end{equation}

La idea principal es aproximar la solución mediante una sucesión de puntos utilizando la pendiente de la función.

\section{Método de Euler}

Dado un paso $h > 0$, el método de Euler se define por:

\begin{equation}
y_{n+1} = y_n + h f(x_n, y_n)
\end{equation}

donde:
\begin{itemize}
\item $x_{n+1} = x_n + h$
\item $y_n$ es la aproximación de la solución en $x_n$
\end{itemize}

\subsection{Interpretación geométrica}

El método utiliza la recta tangente en el punto $(x_n, y_n)$ para aproximar el valor de la solución en el siguiente punto.

\section{Ejemplo Lineal}

Consideremos el problema:

\begin{equation}
\frac{dy}{dx} = y, \quad y(0) = 1
\end{equation}

Sabemos que la solución exacta es:

\begin{equation}
y(x) = e^x
\end{equation}

Aplicamos Euler con paso $h = 0.1$.

\subsection*{Cálculo}

\begin{align*}
y_1 &= y_0 + h f(x_0, y_0) = 1 + 0.1(1) = 1.1 \\
y_2 &= 1.1 + 0.1(1.1) = 1.21 \\
y_3 &= 1.21 + 0.1(1.21) = 1.331
\end{align*}

\section{Ejemplo No Lineal}

Consideremos:

\begin{equation}
\frac{dy}{dx} = x^2 + y, \quad y(0) = 1
\end{equation}

\subsection*{Cálculo con $h = 0.1$}

\begin{align*}
y_1 &= 1 + 0.1(0^2 + 1) = 1.1 \\
y_2 &= 1.1 + 0.1(0.1^2 + 1.1) = 1.1 + 0.1(1.11) = 1.211 \\
y_3 &= 1.211 + 0.1(0.2^2 + 1.211) = 1.211 + 0.1(1.251) = 1.3361
\end{align*}

\section{Análisis de Error}

\subsection{Error Local}

El error local de truncamiento mide el error cometido en un solo paso, suponiendo que el valor inicial es exacto:

\begin{equation}
\tau_n = \frac{y(x_{n+1}) - y(x_n)}{h} - f(x_n, y(x_n))
\end{equation}

Para el método de Euler:

\begin{equation}
\tau_n = O(h)
\end{equation}

\subsection{Error Global}

El error global es la diferencia entre la solución exacta y la aproximación numérica:

\begin{equation}
e_n = y(x_n) - y_n
\end{equation}

Para Euler:

\begin{equation}
e_n = O(h)
\end{equation}

\subsection{Consistencia}

Un método es consistente si el error local tiende a cero cuando $h \to 0$:

\begin{equation}
\lim_{h \to 0} \tau_n = 0
\end{equation}

El método de Euler es consistente porque su error local es $O(h)$.

\subsection{Estabilidad}

La estabilidad estudia cómo los errores se propagan a lo largo de los pasos.

Consideremos la ecuación test:

\begin{equation}
\frac{dy}{dx} = \lambda y
\end{equation}

El método de Euler da:

\begin{equation}
y_{n+1} = (1 + h\lambda) y_n
\end{equation}

Para que el método sea estable, se requiere:

\begin{equation}
|1 + h\lambda| < 1
\end{equation}

\subsection{Precisión}

El método de Euler es de \textbf{orden 1}, lo que significa que el error global es proporcional a $h$:

\begin{equation}
\text{Error global} = O(h)
\end{equation}

Esto implica que para mejorar la precisión es necesario disminuir el tamaño del paso.



\end{document}